\section{Introduction}

Biosafety guidance in the United States and internationally places the obligation to
assess risk on the individual laboratory. The sixth edition of \emph{Biosafety in
Microbiological and Biomedical Laboratories} and the fourth edition of the WHO
\emph{Laboratory Biosafety Manual} both frame risk assessment as a local, procedure-aware
activity rather than a matter of agent classification alone~\cite{bmbl6,wholbm4}. ISO
35001:2019 formalises the same expectation as a management-system requirement, obliging
organisations to identify hazards, assess risk, and mitigate it as a documented
cycle~\cite{iso35001}.

The obligation is clear. The method is not. A laboratory director who wishes to know
whether uncapping evacuated tubes at the open bench is acceptable, or whether colony
picking for mass-spectrometric identification requires a biological safety cabinet, has no
published instrument that operates at that resolution. Agent-based risk group
classifications describe organisms, not procedures. Facility-level frameworks describe
containment levels, not bench steps. The gap between them is where the daily work of a
diagnostic laboratory actually happens, and it is filled, in practice, by local custom.

The scale of the gap is substantial. The Centers for Medicare \& Medicaid Services record
301{,}284 laboratories registered under the Clinical Laboratory Improvement Amendments as
of September 2026, of which 15{,}214 hold a certificate of compliance and 15{,}542 a
certificate of accreditation~\cite{cms2026}. Those 30{,}756 non-waived settings perform
the moderate- and high-complexity testing in which the bench steps considered here are
routine. Each is required to assess its own procedures, and each begins from a blank form.

A second problem compounds the first, and is less often stated. Where procedure-level
guidance does exist, it rarely distinguishes recommendations grounded in measurement from
recommendations grounded in precedent. The recent review by Ravnholdt and
colleagues~\cite{ravnholdt2025} makes this explicit for aerosol-generating procedures,
observing that while some guidance is informed by empirical data, a significant gap
remains in the empirical evidence supporting it, and that highly automated laboratory
equipment in particular has not been characterised. A practitioner reading a control
recommendation cannot tell whether it rests on an air-sampling study, on a cluster of
fatal infections, or on tradition. These three warrant very different degrees of
confidence, and conflating them is not a neutral simplification: it makes the evidence
base appear uniform when it is not.

This paper addresses both problems with one instrument. I report:

\begin{enumerate}[label=(\roman*)]
\item an open, procedure-level scoring instrument of twelve anchored items across three
      dimensions, with an evidence grade reported alongside every score and deliberately
      never combined into it;
\item a mapping of every instrument element to ISO 35001:2019 clause numbers and titles,
      so that a score can be tied to the management-system requirement it serves;
\item reference scores, with cited evidence and recommended control sets, for seven
      routine bench steps under three agent profiles; and
\item an evidence map that makes the strength of the underlying literature visible
      step by step, including where it is absent.
\end{enumerate}

\noindent The principal empirical result is an inversion, described in
Section~\ref{sec:inversion}: the bench steps carrying the highest residual biorisk in this
reference are those for which no aerosol measurement exists, while the steps most often
named in requests for containment have been measured and found to release nothing
detectable. I argue in Section~\ref{sec:why} that this is a property of how the evidence
base was assembled rather than of the hazards themselves.

\section{Methods}

\subsection{Scope and unit of analysis}

The unit of analysis is the bench step: a bounded sequence of manipulations performed on a
specimen or isolate, with a defined start and end, that a laboratory would recognise as
one task. Seven steps were selected to span the routine diagnostic workflow from specimen
arrival to organism identification (Table~\ref{tab:steps}). Each step carries an explicit
inclusion and exclusion boundary in the released data, so that a laboratory applying the
instrument can tell whether its own practice falls inside the step as scored.

Scope is limited to laboratories performing moderate- and high-complexity testing under
CLIA. Research laboratories, high-containment facilities, and veterinary and food testing
are out of scope, as are the pre-laboratory phases of collection and transport.

\subsection{Design principle: score and evidence reported separately}

The instrument's governing design choice is that the score for a step and the strength of
the evidence behind it are reported together and never combined into a single figure.

The rationale is that these two quantities answer different questions. A score answers
``how much residual biorisk does this step carry?''. An evidence grade answers ``how much
should you trust that number?''. A step may legitimately carry a high score on weak
evidence, which is a call for caution and for research; or a low score on strong evidence,
which is a basis for reallocating resources. Averaging the two conflates a statement about
the world with a statement about our knowledge of it, and destroys precisely the
information a frontline laboratory needs in order to decide whether to act now or to
investigate further. Conventional risk matrices, which return a single cell, cannot
express this distinction.

\subsection{Dimension A: inherent aerosol release}

Dimension A scores the capacity of the step to put respirable material into the air,
\emph{before any control is credited}. It comprises five items, each anchored at five
levels: A1 energy applied to the specimen, A2 containment of the specimen during the
step, A3 volume and concentration of biological material, A4 proximity to the operator
breathing zone, and A5 throughput and repetition.

Two features warrant comment. First, A5 is included because a small per-event release
repeated several hundred times per shift is not a small exposure, and procedure-level
assessments that score a single instance of a step systematically understate
high-throughput work. A5 is scored per operator per shift.

Second, A2 is \emph{reverse-coded}. Its anchors run from 1 (fully open; specimen exposed
to room air throughout) to 5 (fully closed; no container opened, pierced, or spun at any
point), which is the direction a practitioner naturally reads containment, but the
opposite of the direction of dimension A as a whole, in which a higher value means more
release. The instrument therefore applies the transformation $A2' = 6 - A2$ within the
aggregation while leaving the published anchor wording unchanged, and the item carries an
explicit direction note. This is recorded here because an earlier draft of this work
aggregated A2 without the transformation, with the consequence that better containment
raised the release score; the defect was detected during verification and is discussed in
Section~\ref{sec:limitations} as a general caution about mixed-direction anchors.

Dimension A is a property of the step. Its mean is the Inherent Aerosol Release Score,
$\IARS \in [1,5]$.

\subsection{Dimension B: scenario consequence}

Dimension B scores what follows if released material reaches a worker. It comprises three
items: B1 inhalational infectivity, B2 severity against the availability of prophylaxis
and treatment (reverse-anchored, so that a high value denotes a bad outcome that is hard
to prevent after exposure), and B3 likelihood of unsuspected presence.

B3 is the item most often absent from procedure-level risk assessments and, on the
evidence reviewed here, the one that matters most. A laboratory applies heightened
precautions only when it knows what it is holding. The organisms that have caused
laboratory-acquired infection in clinical settings are overwhelmingly those that arrived
without warning in a routine specimen and were worked up as ordinary
isolates~\cite{ackelsberg2020,sejvar2005}.

Because dimension B is a property of the organism rather than of the step, it is held as a
scenario parameter. Three agent profiles are defined and reported side by side rather than
collapsing to a single assumed organism (Table~\ref{tab:profiles}). Each step is also
assigned a \emph{primary profile}, being the realistic operating case for that step; where
a step exists only because a specific hazard is suspected, as with cartridge-based
tuberculosis testing, only the applicable profile is reported. The mean of the three items
is the Scenario Consequence Score, $\SCS \in [1,5]$.

\subsection{Dimension C: control reliability}

Dimension C scores how far the controls in place can be relied upon, across four items:
C1 engineering control effectiveness for the release mode of this specific step, C2
independence from human performance, C3 verification, certification and validation, and
C4 personal protective equipment adequacy \emph{for the inhalation route}.

C2 encodes the hierarchy-of-controls intuition in scorable form: a control that works only
if a tired operator remembers it at the end of a night shift is weaker than one built into
the consumable. C4 is deliberately strict. A laboratory coat and gloves protect against
contact and splash, not inhalation, and a step scored on inhalation risk earns no credit
for them.

Dimension C is scored twice for every step: once for \emph{baseline practice} as commonly
described in the literature and guidance, and once for the \emph{recommended control set}
derived in this work. The difference between the two states is what the recommendation is
worth. The mean is the Control Reliability Score, $\CRS \in [1,5]$, in which a higher value
denotes more reliable control.

\subsection{Composition}

The three dimensions compose into a Residual Biorisk Index:
\begin{equation}
\RBI \;=\; \max\!\bigl(1,\; \IARS - \kappa(\CRS - 1)\bigr)\times \SCS ,
\label{eq:rbi}
\end{equation}
with control credit $\kappa = 0.8$ and $\RBI \in [1,25]$. Bands are drawn at 4, 9, 14 and
19, giving Low, Moderate, Elevated, High and Very high.

Two properties of Eq.~\eqref{eq:rbi} are intentional. The release term is floored at 1, so
that no control set is credited with eliminating release entirely; the evidence base
reviewed here does not support such a claim for any of the seven steps. And a control set
at the top of every scale ($\CRS = 5$) reduces the release term by $4\kappa = 3.2$ points,
which can take the maximum release score of 5 down to 1.8 but never to zero.

Both $\kappa$ and the band cut points are set by the author and are asserted rather than
derived. Sensitivity of the step ordering to these choices has not been tested and is
identified in Section~\ref{sec:limitations} as a priority for the next release.

\subsection{Evidence grading}

Every step receives an evidence grade on a five-level scale, assigned from the extracted
evidence rather than from the rater's impression:

\begin{itemize}
\item[\textbf{E1}] Direct quantitative measurement of aerosol release during this exact
      step, in a clinical or realistically simulated clinical setting.
\item[\textbf{E2}] Quantitative measurement of a closely analogous step, or of this step
      in a non-clinical setting, requiring a stated inference to transfer.
\item[\textbf{E3}] Epidemiological or incident evidence linking the step to
      laboratory-acquired infection or documented exposure, without aerosol measurement.
\item[\textbf{E4}] Authoritative guidance or expert consensus only.
\item[\textbf{E5}] No evidence located.
\end{itemize}

A \emph{conditionality flag} is set where the evidence splits sharply between the step
performed as intended and the step performed with a named failure present
(\texttt{failure-conditional}), or where it turns on operator technique
(\texttt{technique-conditional}). Where the flag is set, both states are reported, because
a single figure would misrepresent both.

\subsection{Mapping to ISO 35001:2019}

Each instrument element and each step is mapped to the ISO 35001:2019 clauses it speaks
to, using clause numbers and titles only. No requirement text, figure, or definition from
the standard is reproduced, and no purchase of the standard is required to use this
reference. Across the instrument and the seven steps, 17 distinct clauses are referenced,
concentrated in 6.1.1--6.1.4 (hazard identification, risk assessment, risk mitigation,
performance evaluation), 8.1, 8.3, 8.6 and 8.7 (operational planning, maintenance and
certification, good microbiological technique, PPE), and 10.2 (incident, nonconformity and
corrective action).

The mapping is structural. It indicates which requirement an element serves. It is not a
conformity assessment, and applying the instrument does not make a laboratory conformant.

\subsection{Evidence identification and extraction}

Sources were identified by targeted searching of PubMed Central, open-access journals, and
the published guidance of CDC, WHO and the American Society for Microbiology, restricted
to material accessible without payment. Fifteen sources were admitted. For each step, I
extracted findings bearing on aerosol release or on documented exposure, recording the
finding, its direction (positive, negative, gap, or control recommendation), and its
source. Eighteen findings were extracted in total.

Three claims identified during extraction were \emph{excluded} because they could be
traced only to secondary description and not read at source: a report that direct smear
with matrix solution alone prevents growth of \emph{Brucella} on a MALDI-TOF target; a
description of a simulated incorrect-use condition in the cartridge-assay bioaerosol
study; and the designation and interval of the recognised standard for biological safety
cabinet field certification. Each is held in a register of pending sources, cited nowhere,
and its exclusion is noted where it affects a control set. The first exclusion makes the
MALDI-TOF control set more cautious than the full evidence may warrant; the third is the
reason this reference requires each laboratory to record the certification standard and
interval it applies rather than naming one.

\subsection{Reproducibility}

The instrument, step definitions, agent profiles, source register, and extracted evidence
are held as structured text. Scores are computed by script; no score is entered by hand
into any table. Every numerical statement in this manuscript is interpolated at build time
from a single machine-readable statistics file written by the analysis pipeline, so that
text and data cannot diverge. A quality-control script independently recomputes every
index from the raw sub-scores, checks referential integrity of every citation and clause
reference, and verifies that no value has been imputed; it reports 21 passing checks and
no failures. A publication gate refuses to release while any draft marker, unresolved
verification tag, or placeholder remains.
